%%%%%
%
\section{Exponentiation}

\itemstyle{Op\'erateur~:} {\sf Pow\_f}
\index{Pow\_f}

%%%%%OBJET%%%%
\itemstyle{Objet~:}
Calcul de $base$ � la puisssance $(X+echelle)$, pour des $X$ de type \texttt{float}.

%%%%%ARGUMENTS%%%%
\bigskip
%%%%%ENTREES%%%%
\begin{tabarg}
\multicolumn{3}{l}{\titretab {Entr\'ees}}\\
\hline
$X$ & Vecteur, {\tt float} & Op�rande \\
$base$ & Coeff., {\tt double} & base\\
$scale$ & Coeff., {\tt double} & facteur d'�chelle (en log.)\\
\end{tabarg}

\medskip
%%%%%SORTIES%%%%
\begin{tabarg}
\multicolumn{3}{l}{\titretab {Sorties}}\\
\hline
$Z$ & Vecteur, {\tt float} & R�sultat\\
\end{tabarg}


%%%%%FORMULE%%%%
\itemstyle{Formule}

%%Pow_f
\[
Z = base^{X+scale}
\]



%%%%%COMMENTAIRE%%%%
\itemstyle{Commentaire~:}

Pas de traitement sp\'ecifique des cas hors du domaine de
d�finition d'une puisssance (Si $base$ est nulle, les
�chantillons de $X$ doivent �tre positifs, et si $base$
est n�gative, les
�chantillons de $X$ doivent �tre entiers).

%%%%%AUTEUR%%%%
\itemstyle{Auteur~:}
CNRM/GMEI/TRAMM

%%%%%REFERENCE%%%%
\itemstyle{R\'ef\'erences~:}
